\section{Introduction}
\label{sec:intro}

\subsection{Research need and gap}

The conceptual design phase establishes the foundation of every
aircraft development program by defining the overall configuration,
mission feasibility, propulsion architecture, structural layout,
subsystem selection, and overall system performance. These early-stage
decisions rely heavily on preliminary sizing calculations, engineering
heuristics, multidisciplinary trade-off analyses, and numerous design
iterations, despite limited available information
\citep{corrado2022,karali2024surrogate}. In recent years, artificial
intelligence (AI) has demonstrated considerable potential in supporting
individual engineering tasks; however, existing applications remain
fragmented and are typically limited to isolated activities such as
surrogate modeling, geometry generation, programming assistance,
optimization, or document preparation
\citep{karali2024mdo,du2026,pliakos2025}. There is still no structured
methodology that systematically integrates AI into the complete
conceptual aircraft design process while preserving physics-based
engineering calculations, multidisciplinary reasoning, and engineering
verification. This need is particularly critical for UAV development,
where rapidly evolving operational requirements, frequent design
updates, and short development cycles demand conceptual design
iterations significantly faster than those traditionally adopted in
conventional aircraft programs \citep{karali2023,shakeri2018}.

A qualitatively new opportunity has emerged with large language model
(LLM) based \emph{agentic} coding tools, which can read, write,
execute, and test an entire engineering codebase under natural-language
direction. Unlike the surrogate-model and optimizer literature, which
replaces expensive physics with learned approximations, an LLM agent
can operate one level higher: it can \emph{build and maintain} the
deterministic physics toolchain itself, propose and implement design
changes across every affected discipline simultaneously, and document
the reasoning --- while the engineering numbers continue to come
exclusively from transparent, testable, physics-based code. How to
harness this capability without sacrificing rigor, traceability, and
verification is, to the authors' knowledge, unaddressed in the aircraft
design literature.

\subsection{Proposed methodology}

This study presents an AI-assisted conceptual aircraft design
methodology that supports the complete early-stage development process
from mission definition to a validated conceptual configuration. The
methodology integrates preliminary aircraft sizing, propulsion sizing,
energy storage estimation, subsystem and equipment selection,
configuration development, structural considerations, and
multidisciplinary engineering evaluations within a unified iterative
framework. Artificial intelligence functions as an engineering
collaborator that accelerates engineering reasoning, alternative
generation, knowledge synthesis, component evaluation, technical
documentation, and design-space exploration, while analytical
calculations remain the basis of all engineering decisions. Numerical
analyses, computational simulations, and external expert findings are
continuously incorporated into successive design iterations, enabling
the conceptual design to evolve through evidence-based refinement
rather than linear development.

The methodology is realized concretely as a \emph{governed
human--AI--repository loop}: the design lives entirely in a
version-controlled repository whose architecture (configuration
$\rightarrow$ physics library $\rightarrow$ notebook pipeline
$\rightarrow$ typed hand-offs $\rightarrow$ CAD/CFD consumers) is
itself part of the method; the AI agent works inside that repository
under a machine-readable set of engineering rules; and every change is
gated by human review, architecture decision records, pinned
design-point regression tests, and continuous integration that
re-executes the full design. To demonstrate applicability, a
tail-sitter VTOL UAV is developed as an illustrative case study,
covering the process from mission requirements to preliminary sizing,
subsystem selection, configuration development, commercial
component freeze, and hand-off to higher-fidelity validation.

The work is guided by the following research question:

\begin{quote}
\textbf{RQ:} \emph{How can AI accelerate conceptual aircraft design
while preserving physics-based engineering rigor?}
\end{quote}

\subsection{Main contributions}

The primary contribution of this work is the introduction of a
structured AI-assisted methodology for conceptual aircraft design that
integrates preliminary engineering calculations, multidisciplinary
subsystem selection, engineering knowledge management, numerical
analyses, and iterative design refinement into a single coherent
workflow. Unlike existing AI applications that support isolated
engineering tasks, the proposed methodology formalizes the interaction
between human engineering expertise, an LLM-based AI agent, analytical
models, and validation activities throughout the conceptual design
process --- including the guardrails (design-point regression pins,
one-way dataflow, decision records, standing-finding discipline) that
make AI acceleration compatible with engineering accountability. The
framework is aircraft-independent and can be adapted to a wide range of
aerospace platforms that require rapid, multidisciplinary conceptual
development.

As a second contribution, the paper demonstrates the practical
implementation of the methodology through the conceptual design of a
tail-sitter VTOL UAV, resulting in a complete preliminary sizing
procedure, propulsion system selection, onboard equipment architecture,
and configuration development, together with quantitative process
evidence extracted from the project's own git history: the share of
AI-co-authored engineering commits, the cadence of design-point
snapshots, and the mechanisms by which errors were caught. Furthermore,
the tail-sitter case study provides application-specific engineering
insights into the conceptual development of this unconventional
aircraft configuration, illustrating how AI-assisted multidisciplinary
workflows can reduce development effort while maintaining engineering
rigor, traceability, and design quality.

The remainder of the paper is organized as follows.
Section~\ref{sec:related} reviews AI applications in conceptual
aircraft and UAV design and motivates the tail-sitter as a validation
case. Section~\ref{sec:method} presents the methodology.
Section~\ref{sec:results} reports the case study results and process
metrics and discusses them against the research question.
Section~\ref{sec:conclusions} concludes.
